
the morphisms @@M0@@ for @@M1@@ (resp. @@M2@@) from the special and generic fibers to the ambient scheme. Denote by
@@M3@@
the induced morphisms, which therefore give rise to the Cartesian squares
\sgaStableTarget{sga6.e10.eq.7-10-0}@@M4@@

First suppose that @@M5@@ is \emph{proper}. I claim that the following diagram then commutes:
\sgaStableTarget{sga6.e10.eq.7-10-1}@@M6@@
This result generalizes 7.5 (since, for @@M7@@, the arrows of (7.3.2) are identified with the identity map of @@M8@@). Note that, when @@M9@@ and @@M10@@ satisfy the conditions of 7.7, i.e., have regular fibers, @@M11@@ and @@M12@@ automatically have finite Tor-dimension, and the diagram can be interpreted as a commutative diagram with @@M13@@, which we leave the reader to rewrite.

To prove the commutativity of (7.10.1), one sees at once, using the definitions of @@M14@@ and @@M15@@, that one is reduced to proving the commutativity of the two front squares of the prism
\sgaStableTarget{sga6.e10.eq.unnamed}@@M16@@
